% 海洋行星（综述）
% license CCBYSA3
% type Wiki

本文根据 CC-BY-SA 协议转载翻译自维基百科\href{https://en.wikipedia.org/wiki/Ocean_world}{相关文章}。

\begin{figure}[ht]
\centering
\includegraphics[width=6cm]{./figures/76246827a4922672.png}
\caption{地球表面主要由海洋覆盖，海洋占据了地球表面的 $75\%$。因此，地球可以被视为一个水世界，但并非一个完全的海洋行星。} \label{fig_Ocwo_1}
\end{figure}
海洋世界（ocean world）、海洋行星（ocean planet）或水世界（water world）是一类行星或天然卫星，其水圈中包含大量以海洋形式存在的水，这些海洋可以位于表面之下，形成地下海洋，或直接存在于表面，甚至可能淹没所有陆地$^{[1][2][3][4]}$。“海洋世界”一词有时也被用于指代其海洋由其他流体或造海物质（thalassogen）构成的天体$^{[5]}$，例如熔岩（如木卫一 Io）、氨（与水形成共晶混合物，可能存在于土卫六 Titan 的内部海洋中），或碳氢化合物（如土卫六表面的情况，其可能是系外海洋中最常见的一类）$^{[6]}$。对地外海洋的研究被称为行星海洋学（planetary oceanography）。

地球是目前唯一已知在其表面存在液态水体的天文天体，尽管在木星的卫星欧罗巴（Europa）与盖尼米德（Ganymede），以及土星的卫星恩克拉多斯（Enceladus）与土卫六（Titan）上，人们推测其表面之下可能存在地下海洋$^{[7]}$。已经发现若干系外行星具备支持液态水存在的适宜条件$^{[8]}$。此外，在地球上也发现了大量地下水，主要以含水层的形式存在$^{[9]}$。对于系外行星而言，受限于当前技术，尚无法直接观测其表面的液态水，因此通常以大气中的水汽作为替代性指标$^{[10]}$。海洋世界的性质为理解其自身历史以及整个太阳系的形成与演化提供了重要线索；与此同时，它们在生命起源与维持方面的潜力也引起了广泛关注。

2020 年 6 月，NASA 科学家基于数学建模研究报告称，在银河系中拥有海洋的系外行星很可能是普遍存在的$^{[11][12]}$。
\subsection{概述}  
\subsubsection{定义}  
根据 Lunine 的定义，“海洋”被界定为“稳定存在、环绕整个天体的液态水体”$^{[13]}$。此外，“海洋世界”这一术语用于指代太阳系中那些拥有稳定、覆盖全球的液态水体的天体；与之相对，“海洋行星”和“水世界”则专指系外行星（即绕其他恒星运行的行星），其整体组成中含有较高质量分数的水$^{[13]}$。
\subsubsection{太阳系行星天体}
\begin{figure}[ht]
\centering
\includegraphics[width=8cm]{./figures/4b85d59dce5a977e.png}
\caption{土卫二（Enceladus）内部结构示意图} \label{fig_Ocwo_2}
\end{figure}
海洋世界因其可能孕育生命并维持生物活动的潜力而受到天体生物学家的高度关注$^{[4][3]}$。太阳系中被认为拥有地下海洋的主要卫星和矮行星尤为重要，因为与远在数光年之外、受当前技术所限难以直接探测的系外行星不同，这些天体可以通过空间探测器抵达并开展实地研究。除地球之外，太阳系中证据最为充分的水世界包括卡利斯托（Callisto）、恩克拉多斯（Enceladus）、欧罗巴（Europa）、盖尼米德（Ganymede）以及土卫六（Titan）$^{[3][14]}$。其中，欧罗巴和恩克拉多斯由于其外壳较薄且存在冰火山活动，被视为极具吸引力的探测目标。

太阳系中的其他天体也被视为可能拥有地下海洋的候选者，这一判断基于单一类型的观测证据或理论建模结果，包括天卫一（Ariel）$^{[14]}$、天卫三（Titania）$^{[15][16]}$、天卫二（Umbriel）$^{[17]}$、谷神星（Ceres）$^{[3]}$、土卫四（Dione）$^{[18]}$、土卫一（Mimas）$^{[19][20]}$、天卫五（Miranda）$^{[14]}$、天卫四（Oberon）$^{[4][21]}$、冥王星（Pluto）$^{[22]}$、海卫一（Triton）$^{[23]}$、厄里斯（Eris）$^{[4][24]}$以及鸟神星（Makemake）$^{[24]}$。
\subsubsection{系外行星}
\begin{figure}[ht]
\centering
\includegraphics[width=8cm]{./figures/d53b32c9e8bb7ae8.png}
\caption{一组含水、大小各异的系外行星，与地球进行对比（艺术家概念图；2018 年 8 月 17 日）$^{[25]}$} \label{fig_Ocwo_3}
\end{figure}
\begin{figure}[ht]
\centering
\includegraphics[width=8cm]{./figures/c7e845b1bfbd19f5.png}
\caption{以纯海洋世界作为过渡群体的系外行星族群示意图：它们位于气态巨行星与熔岩行星或岩石行星之间。} \label{fig_Ocwo_4}
\end{figure}
在太阳系之外，被描述为**候选海洋世界**的系外行星包括 GJ 1214 b $^{[26][27]}$、Kepler-22b、Kepler-62e、Kepler-62f $^{[28][29][30][31]}$，以及 Kepler-11 $^{[32]}$ 和 TRAPPIST-1 $^{[33][34]}$ 行星系统中的行星。

近年来，系外行星 TOI-1452 b、Kepler-138c 和 Kepler-138d 被发现其密度与其质量中有很大一部分由水构成的情形相一致 $^{[35][36]}$。此外，对高质量岩石行星 LHS 1140 b 的模型研究表明，其表面可能被一片深海所覆盖 $^{[37]}$。

尽管地球表面有 70.8\% 被水覆盖 $^{[38]}$，水在地球总质量中仅占约 0.05\%。在地外海洋世界中，海洋可能深且致密，以至于即便在高温条件下，巨大的压力也会使水转变为冰。在这类海洋的下部区域，数千巴的巨大压力可能促成奇异形态冰（如冰 V）地幔的形成 $^{[32]}$。这种冰并不一定像通常意义上的冰那样寒冷。若行星距离其恒星足够近，使水达到沸点，则水将进入超临界状态，不再具有清晰界定的表面 $^{[39]}$。即便是在较冷、以水为主的行星上，其大气层也可能比地球厚得多，并主要由水汽构成，从而产生极强的温室效应。此类行星必须足够小，无法保留厚重的氢和氦包层 $^{[40]}$，或者足够靠近其主星，以致这些轻元素被剥离 $^{[32]}$；否则，它们将演化成类似天王星和海王星那样的、更温暖版本的冰巨行星（此处尚需进一步证据）。
\subsection{历史}
20 世纪初的引力计算曾提出，木卫二（Europa）的组成可能富含水。1957 年，杰拉德·柯伊伯（Gerard Kuiper）通过地基观测证实了其表面主要由水冰构成。$^{[41]}$

在 20 世纪 70 年代行星探测任务之前，相关的重要理论基础工作已相继完成。1971 年，Lewis 指出，仅放射性衰变就可能足以在大型卫星中形成地下海洋，尤其是在存在氨（NH$_3$）的情况下。1979 年，Peale 和 Cassen 揭示了潮汐加热（亦称潮汐形变）在卫星演化与内部结构中的关键作用。$^{[3]}$
1992 年，人类首次确认探测到一颗系外行星。2003 年，Marc Kuchner 以及 2004 年 Alain Léger 等人提出，一小部分在雪线之外形成的冰质行星可能向内迁移至约 $1,\mathrm{AU}$ 的位置，其外层随后发生融化。$^{[42][43]}$

来自哈勃空间望远镜，以及先锋号、伽利略号、旅行者号、卡西尼–惠更斯号和新视野号等任务所积累的证据，强烈表明太阳系外侧的多个天体在隔热的冰壳之下蕴藏着内部液态水海洋。$^{[3][44]}$
与此同时，于 2009 年 3 月 7 日发射的开普勒空间望远镜已发现数千颗系外行星，其中约有 50 颗大小与地球相当，位于宜居带内或附近。$^{[45][46]}$

目前已探测到质量、尺寸和轨道各异的行星，这不仅展示了行星形成过程的多样性，也反映了行星在形成之后通过环恒星盘发生迁移的普遍现象。$^{[10]}$
截至 2025 年 10 月 30 日，已确认的系外行星共有 6,128 颗，分布在 4,584 个行星系统中，其中 1,017 个系统包含多于一颗行星。$^{[47]}$

2020 年 6 月，NASA 科学家基于数学建模研究报告称，银河系中拥有海洋的系外行星很可能相当普遍。$^{[11]}$

2022 年 8 月，距离地球约 99 光年的超级地球系外行星 TOI-1452 b 被凌日系外行星巡天卫星发现，该行星可能拥有深层海洋。$^{[35]}$
\subsection{形成}
\begin{figure}[ht]
\centering
\includegraphics[width=6cm]{./figures/00fa51d11bed7959.png}
\caption{阿塔卡马大型毫米/亚毫米阵列（ALMA）拍摄的 HL 金牛座（HL Tauri）原行星盘图像} \label{fig_Ocwo_5}
\end{figure}